% 负责人：A / hycx233。这里只汇总已保存结果，最终定稿随 B 的评价更新。
\begin{abstract}
本项目使用大肠杆菌 WT、$\Delta stpA$ 和 $\Delta hns\Delta stpA$ 三种条件、
每种两个重复的 Micro-C 数据，完成已知结构分类、候选发现、多轨道可视化及条件差异分析。
将 10 bp 接触矩阵稀疏聚合为 100 bp 工作数据，统一坐标、环状窗口、有效掩码与空间分组，
分别保留深度归一化强度和 O/E 形态输入。344 条已知结构形成 688 个 WT 窗口，
B 按固定参数独立复跑的轻量 CNN 在 104 个测试窗口上的准确率为 0.5096，
Macro-F1 为 0.3454；逻辑回归基线对应为 0.8365 和 0.5644，当前 CNN 未超过该基线。
CNN、逻辑回归和多数类基线均未正确识别测试集中的 8 个 CHID 窗口，说明少数类识别仍有限。
自编码器扫描 4,642 个窗口得到 47 个高分候选，合并后为 10 个独立位置；
结合已知参考与聚类规则，未发现满足本组疑似新簇条件的结果。
三条件轨道共输出 465 张图。对 344 条已知结构的 688 次条件比较中，
$\Delta stpA$ 和双缺失条件分别有 217、182 条结构满足预设的一致增强描述规则。
这些变化反映测序深度归一化后的相对强度，不作为统计显著性或单因子因果证据。
WT 跨重复分析将 47 个候选窗口合为 10 个独立位置，并对每个位置预先固定的一个代表窗口进行比较；Pearson 相关中位数为 0.8572。该结果仅支持共同有效区域的图形一致性，不等同于逐一验证全部候选窗口；按既定规则，本次可信新结构为 0。
\end{abstract}
